\usepackage{xstring}
% cmap must precede fontenc so it can wrap the CMap entries that fontenc
% pulls in; otherwise the package issues a runtime warning and some
% glyphs are left unprocessed in the resulting PDF.
\usepackage{cmap}
\usepackage[T1, T2A]{fontenc}
\usepackage[utf8]{inputenc}
\usepackage[english]{babel}
\usepackage{textgreek}
\usepackage{pmboxdraw}
\usepackage{newunicodechar}

% Box-drawing glyphs (pmboxdraw provides ASCII fallbacks for these).
% See https://tex.org.uk/macros/latex/contrib/pmboxdraw/pmboxdraw.pdf
\newunicodechar{╔}{\textSFxxxix}
\newunicodechar{═}{\textSFxliii}
\newunicodechar{╗}{\textSFxxv}
\newunicodechar{║}{\textSFxxiv}
\newunicodechar{╚}{\textSFxxxviii}
\newunicodechar{╝}{\textSFxxvi}
\newunicodechar{╭}{\textSFi}
\newunicodechar{─}{\textSFx}
\newunicodechar{╮}{\textSFiii}
\newunicodechar{│}{\textSFxi}
\newunicodechar{╰}{\textSFii}
\newunicodechar{╯}{\textSFiv}

% Other Unicode glyphs that appear verbatim in source comments and the
% project's docs. pdftex with [utf8]{inputenc} only knows a small core
% set; everything else has to be mapped explicitly. Math glyphs go
% through ensuremath so they render the same in text and math mode.
\newunicodechar{×}{\ensuremath{\times}}
\newunicodechar{–}{\textendash}
\newunicodechar{—}{\textemdash}
\newunicodechar{…}{\textellipsis}
\newunicodechar{←}{\ensuremath{\leftarrow}}
\newunicodechar{→}{\ensuremath{\rightarrow}}
\newunicodechar{≤}{\ensuremath{\leq}}
\newunicodechar{≥}{\ensuremath{\geq}}
\newunicodechar{█}{\rule{0.6em}{0.7em}}      % full block
\newunicodechar{░}{\rule{0.6em}{0.7em}}      % light shade — best-effort
\newunicodechar{▶}{\ensuremath{\blacktriangleright}}
\newunicodechar{✓}{\checkmark}

\renewcommand{\sfdefault}{cmss}
\renewcommand{\ttdefault}{cmtt}
\renewcommand{\rmdefault}{cmss}

% Font for footnotes.
\renewcommand{\footnotesize}{\fontsize{12}{13}\selectfont}

\let\footnotesize\scriptsize
